
\documentclass{article}
\usepackage{colortbl}
\usepackage{makecell}
\usepackage{multirow}
\usepackage{supertabular}

\begin{document}

\newcounter{utterance}

\centering \large Interaction Transcript for game `wordle', experiment `medium\_frequency\_words\_no\_clue\_no\_critic', episode 16 with glm{-}4.6{-}t1.0.
\vspace{24pt}

{ \footnotesize  \setcounter{utterance}{1}
\setlength{\tabcolsep}{0pt}
\begin{supertabular}{c@{$\;$}|p{.15\linewidth}@{}p{.15\linewidth}p{.15\linewidth}p{.15\linewidth}p{.15\linewidth}p{.15\linewidth}}
    \# & \multicolumn{2}{c}{Player} && \multicolumn{2}{c}{Game Master} \\
    \hline

    \theutterance \stepcounter{utterance}  
    & & & \multicolumn{4}{p{0.6\linewidth}}{
        \cellcolor[rgb]{0.9,0.9,0.9}{
            \makecell[{{p{\linewidth}}}]{
                \texttt{\tiny{[P1$\langle$GM]}}
                \texttt{You are a language wizard who likes to guess words by using the given rules.} \\
\\ 
\texttt{Welcome to Wordle! You have six attempts to guess the target word, a valid English word of five lowercase letters (a{-}z). Please use the tags "explanation:" and "guess:" to provide a concise explanation for each guess.} \\
\\ 
\texttt{For instance, if your guess is "apple", your response should be} \\
\texttt{explanation: this is a common five{-}letter English word, and I am starting my guess with this word.} \\
\texttt{guess: apple} \\
\\ 
\texttt{After each guess, your answer will be validated, and you will receive feedback indicating which letters are correct (green), which letters are correct but in the wrong position (yellow), and which letters are incorrect (red). This feedback can be useful in determining which letters to include or exclude in your next guess.} \\
\\ 
\texttt{For example, the feedback for "apple" might be:} \\
\texttt{guess\_feedback: a<yellow> p<yellow> p<green> l<yellow> e<red>} \\
\\ 
\texttt{The explanation should contain details about how the guess\_feedback is used to arrive at a new guess.} \\
\\ 
\texttt{Do not generate any other text.} \\
\\ 
\texttt{Let's begin with your first guess.} \\
            }
        }
    }
     \\ \\

    \theutterance \stepcounter{utterance}  
    & \multicolumn{4}{p{0.6\linewidth}}{
        \cellcolor[rgb]{0.561,0.737,0.561}{
            \makecell[{{p{\linewidth}}}]{
                \texttt{\tiny{[P1$\rangle$GM]}}
                \texttt{explanation: I will start with a common word containing a mix of vowels and consonants to gather initial information.} \\
\texttt{guess: crane} \\
            }
        }
    }
    & & \\ \\

    \theutterance \stepcounter{utterance}  
    & & & \multicolumn{2}{p{0.3\linewidth}}{
        \cellcolor[rgb]{0.9,0.9,0.9}{
            \makecell[{{p{\linewidth}}}]{
                \texttt{\tiny{[GM$|$GM]}}
                \texttt{attempts = 1} \\
\texttt{target\_word = canny} \\
\texttt{guess = crane} \\
\texttt{guess\_feedback = c<green> r<red> a<yellow> n<green> e<red>} \\
            }
        }
    }
    & & \\ \\

    \theutterance \stepcounter{utterance}  
    & & & \multicolumn{4}{p{0.6\linewidth}}{
        \cellcolor[rgb]{0.9,0.9,0.9}{
            \makecell[{{p{\linewidth}}}]{
                \texttt{\tiny{[P1$\langle$GM]}}
                \texttt{guess\_feedback: c<green> r<red> a<yellow> n<green> e<red>} \\
\\ 
\texttt{Provide your response only in this format.} \\
\texttt{explanation: details} \\
\texttt{guess: word} \\
            }
        }
    }
     \\ \\

    \theutterance \stepcounter{utterance}  
    & \multicolumn{4}{p{0.6\linewidth}}{
        \cellcolor[rgb]{0.561,0.737,0.561}{
            \makecell[{{p{\linewidth}}}]{
                \texttt{\tiny{[P1$\rangle$GM]}}
                \texttt{explanation: The feedback from 'crane' indicates 'c' and 'n' are in the correct positions, while 'a' is in the word but not in the second position. The letters 'r' and 'e' are excluded. My guess 'china' uses the green letters correctly, tests the yellow 'a' in a new possible position, and introduces two new common letters, 'h' and 'i'.} \\
\texttt{guess: china} \\
            }
        }
    }
    & & \\ \\

    \theutterance \stepcounter{utterance}  
    & & & \multicolumn{2}{p{0.3\linewidth}}{
        \cellcolor[rgb]{0.9,0.9,0.9}{
            \makecell[{{p{\linewidth}}}]{
                \texttt{\tiny{[GM$|$GM]}}
                \texttt{attempts = 2} \\
\texttt{target\_word = canny} \\
\texttt{guess = china} \\
\texttt{guess\_feedback = c<green> h<red> i<red> n<green> a<yellow>} \\
            }
        }
    }
    & & \\ \\

    \theutterance \stepcounter{utterance}  
    & & & \multicolumn{4}{p{0.6\linewidth}}{
        \cellcolor[rgb]{0.9,0.9,0.9}{
            \makecell[{{p{\linewidth}}}]{
                \texttt{\tiny{[P1$\langle$GM]}}
                \texttt{guess\_feedback: c<green> h<red> i<red> n<green> a<yellow>} \\
\\ 
\texttt{Provide your response only in this format.} \\
\texttt{explanation: details} \\
\texttt{guess: word} \\
            }
        }
    }
     \\ \\

    \theutterance \stepcounter{utterance}  
    & \multicolumn{4}{p{0.6\linewidth}}{
        \cellcolor[rgb]{0.561,0.737,0.561}{
            \makecell[{{p{\linewidth}}}]{
                \texttt{\tiny{[P1$\rangle$GM]}}
                \texttt{explanation: With 'c' and 'n' confirmed in their positions, the word pattern is '\_c\_n\_'. 'a' is in the word but not in the third or fifth spot. Placing 'a' in the first position gives the pattern 'ac\_n\_'. The letters 'r', 'e', 'h', and 'i' are excluded. 'acorn' fits all these criteria and is a common English word.} \\
\texttt{guess: acorn} \\
            }
        }
    }
    & & \\ \\

    \theutterance \stepcounter{utterance}  
    & & & \multicolumn{2}{p{0.3\linewidth}}{
        \cellcolor[rgb]{0.9,0.9,0.9}{
            \makecell[{{p{\linewidth}}}]{
                \texttt{\tiny{[GM$|$GM]}}
                \texttt{attempts = 3} \\
\texttt{target\_word = canny} \\
\texttt{guess = acorn} \\
\texttt{guess\_feedback = a<yellow> c<yellow> o<red> r<red> n<yellow>} \\
            }
        }
    }
    & & \\ \\

    \theutterance \stepcounter{utterance}  
    & & & \multicolumn{4}{p{0.6\linewidth}}{
        \cellcolor[rgb]{0.9,0.9,0.9}{
            \makecell[{{p{\linewidth}}}]{
                \texttt{\tiny{[P1$\langle$GM]}}
                \texttt{guess\_feedback: a<yellow> c<yellow> o<red> r<red> n<yellow>} \\
\\ 
\texttt{Provide your response only in this format.} \\
\texttt{explanation: details} \\
\texttt{guess: word} \\
            }
        }
    }
     \\ \\

    \theutterance \stepcounter{utterance}  
    & \multicolumn{4}{p{0.6\linewidth}}{
        \cellcolor[rgb]{0.561,0.737,0.561}{
            \makecell[{{p{\linewidth}}}]{
                \texttt{\tiny{[P1$\rangle$GM]}}
                \texttt{explanation: While the feedback for 'acorn' seemed to conflictingly re{-}mark the previously confirmed letters 'c' and 'n', it did confirm that 'o' and 'r' are not in the word. Sticking with the established positions of 'c' at the start and 'a' being in the word (but not at the end), I will form the 'c\_a\_' pattern. 'l' and 'k' are common consonants. 'clank' fits the known constraints and is a probable target word.} \\
\texttt{guess: clank} \\
            }
        }
    }
    & & \\ \\

    \theutterance \stepcounter{utterance}  
    & & & \multicolumn{2}{p{0.3\linewidth}}{
        \cellcolor[rgb]{0.9,0.9,0.9}{
            \makecell[{{p{\linewidth}}}]{
                \texttt{\tiny{[GM$|$GM]}}
                \texttt{attempts = 4} \\
\texttt{target\_word = canny} \\
\texttt{guess = clank} \\
\texttt{guess\_feedback = c<green> l<red> a<yellow> n<green> k<red>} \\
            }
        }
    }
    & & \\ \\

    \theutterance \stepcounter{utterance}  
    & & & \multicolumn{4}{p{0.6\linewidth}}{
        \cellcolor[rgb]{0.9,0.9,0.9}{
            \makecell[{{p{\linewidth}}}]{
                \texttt{\tiny{[P1$\langle$GM]}}
                \texttt{guess\_feedback: c<green> l<red> a<yellow> n<green> k<red>} \\
\\ 
\texttt{Provide your response only in this format.} \\
\texttt{explanation: details} \\
\texttt{guess: word} \\
            }
        }
    }
     \\ \\

    \theutterance \stepcounter{utterance}  
    & \multicolumn{4}{p{0.6\linewidth}}{
        \cellcolor[rgb]{0.561,0.737,0.561}{
            \makecell[{{p{\linewidth}}}]{
                \texttt{\tiny{[P1$\rangle$GM]}}
                \texttt{explanation: The feedback from the last four guesses consistently confirms 'c' at the second position and 'n' at the fourth position, resulting in the pattern `\_c\_n\_`. The letter 'a' has consistently received yellow feedback, but its position remains contradictory, as all five slots seem to be eliminated. I will hypothesize that the yellow feedback for 'a' in the first guess, 'crane', was an exception. Placing 'a' in the third position gives the pattern `\_can\_`. I will now test this structure with a common combination of remaining letters, 's' and 'd'.} \\
\texttt{guess: scand} \\
            }
        }
    }
    & & \\ \\

    \theutterance \stepcounter{utterance}  
    & & & \multicolumn{2}{p{0.3\linewidth}}{
        \cellcolor[rgb]{0.9,0.9,0.9}{
            \makecell[{{p{\linewidth}}}]{
                \texttt{\tiny{[GM$|$GM]}}
                \texttt{attempts = 5} \\
\texttt{target\_word = canny} \\
\texttt{guess = scand} \\
\texttt{guess\_feedback = s<red> c<yellow> a<yellow> n<green> d<red>} \\
            }
        }
    }
    & & \\ \\

    \theutterance \stepcounter{utterance}  
    & & & \multicolumn{4}{p{0.6\linewidth}}{
        \cellcolor[rgb]{0.9,0.9,0.9}{
            \makecell[{{p{\linewidth}}}]{
                \texttt{\tiny{[P1$\langle$GM]}}
                \texttt{guess\_feedback: s<red> c<yellow> a<yellow> n<green> d<red>} \\
\\ 
\texttt{Provide your response only in this format.} \\
\texttt{explanation: details} \\
\texttt{guess: word} \\
            }
        }
    }
     \\ \\

    \theutterance \stepcounter{utterance}  
    & \multicolumn{4}{p{0.6\linewidth}}{
        \cellcolor[rgb]{0.561,0.737,0.561}{
            \makecell[{{p{\linewidth}}}]{
                \texttt{\tiny{[P1$\rangle$GM]}}
                \texttt{explanation: From the previous five guesses, 'n' is consistently green at position 4. The feedback from 'scand' confirmed that 'a' and 'c' are in the word, but not at positions 3 and 2 respectively. To resolve the contradictions, 'a' must be at position 1 and 'c' must be at position 3. To complete the word, I will introduce two new, common letters that have not been tested: 'g' and 't'. This forms the word 'agent'.} \\
\texttt{guess: agent} \\
            }
        }
    }
    & & \\ \\

    \theutterance \stepcounter{utterance}  
    & & & \multicolumn{2}{p{0.3\linewidth}}{
        \cellcolor[rgb]{0.9,0.9,0.9}{
            \makecell[{{p{\linewidth}}}]{
                \texttt{\tiny{[GM$|$GM]}}
                \texttt{attempts = 6} \\
\texttt{target\_word = canny} \\
\texttt{guess = agent} \\
\texttt{guess\_feedback = a<yellow> g<red> e<red> n<green> t<red>} \\
            }
        }
    }
    & & \\ \\

    \theutterance \stepcounter{utterance}  
    & & & \multicolumn{2}{p{0.3\linewidth}}{
        \cellcolor[rgb]{0.9,0.9,0.9}{
            \makecell[{{p{\linewidth}}}]{
                \texttt{\tiny{[GM$|$GM]}}
                \texttt{game\_result = LOSS} \\
            }
        }
    }
    & & \\ \\

\end{supertabular}
}

\end{document}
